%-----------------------------------
% Define document and include general packages
%-----------------------------------
% Tabellen- und Abbildungsverzeichnis stehen normalerweise nicht im
% Inhaltsverzeichnis. Gleiches gilt für das Abkürzungsverzeichnis (siehe unten).
% Manche Dozenten bemängeln das. Die Optionen 'listof=totoc,bibliography=totoc'
% geben das Tabellen- und Abbildungsverzeichnis im Inhaltsverzeichnis (toc=Table
% of Content) aus.
% Da es aber verschiedene Regelungen je nach Dozent geben kann, werden hier
% beide Varianten dargestellt.
\documentclass[12pt,oneside,titlepage,listof=totoc,bibliography=totoc]{scrartcl}
%\documentclass[12pt,oneside,titlepage]{scrartcl}

%-----------------------------------
% Dokumentensprache
%-----------------------------------
%\def\FOMEN{}% Auskommentieren um die Dokumentensprache auf englisch zu ändern
\newif\ifde
\newif\ifen

%-----------------------------------
% Meta informationen
%-----------------------------------
%-----------------------------------
% Meta Informationen zur Arbeit
%-----------------------------------

% Autor
\newcommand{\myAutor}{Max Mustermann}
\newcommand{\myVorname}{Max}
\newcommand{\myNachname}{Mustermann}

% Adresse
\newcommand{\myAdresse}{Heidestra\ss e 17 \\ \> \> \> 51147 Köln}

% Titel der Arbeit
\newcommand{\myTitel}{LATEX-Vorlage - mit Biblatex}

% Betreuer
\newcommand{\myBetreuer}{Prof. Dr. Peter Lustig}

% Lehrveranstaltung
\newcommand{\myLehrveranstaltung}{Modul Nr. 1}

% Matrikelnummer
\newcommand{\myMatrikelNr}{123456}

% Ort
\newcommand{\myOrt}{Düsseldorf}

% Datum der Abgabe
\newcommand{\myAbgabeDatum}{\today}

% Semesterzahl
\newcommand{\mySemesterZahl}{7}

% Name der Hochschule
\newcommand{\myHochschulName}{FOM Hochschule für Oekonomie \& Management}

% Standort der Hochschule
\newcommand{\myHochschulStandort}{Düsseldorf}

% Studiengang
\newcommand{\myStudiengang}{Wirtschaftsinformatik}

% Art der Arbeit
\newcommand{\myThesisArt}{Bachelor Thesis}

% Zu erlangender akademische Grad
\newcommand{\myAkademischerGrad}{Bachelor of Science (B.Sc.)}

% Firma
\newcommand{\myFirma}{Mustermann GmbH}


\ifdefined\FOMEN
    %Englisch
    \entrue
    \usepackage[english]{babel}
\else
    %Deutsch
    \detrue
    \usepackage[ngerman]{babel}
\fi


\newcommand{\langde}[1]{%
   \ifde\selectlanguage{ngerman}#1\fi}
\newcommand{\langen}[1]{%
   \ifen\selectlanguage{english}#1\fi}
\usepackage[utf8]{luainputenc}
\langde{\usepackage[babel,german=quotes]{csquotes}}
\langen{\usepackage[babel,english=british]{csquotes}}
\usepackage[T1]{fontenc}
\usepackage{fancyhdr}
\usepackage{fancybox}

% SEITENFORMAT %% 
% header/footer = 140mm 
%  Oben: 4 cm (sowie einen Abstand vom Seitenrand der Kopfzeile: 2 cm),
%  Unten: 2 cm (sowie einen Abstand vom Seitenrand der Fußzeile: 2 cm),
% Links: 4 cm,
%  Rechts: 2 cm.

\usepackage[a4paper, left=4cm, right=2cm, top=3.8cm, bottom=3cm,footskip=1cm,headsep=1cm]{geometry}

% weitere Pakete
\usepackage{graphicx}
\usepackage{colortbl}
\usepackage[capposition=top]{floatrow}
\usepackage{array}
\usepackage{float}      %Positionierung von Abb. und Tabellen mit [H] erzwingen
\usepackage{footnote}

% Darstellung der Beschriftung von Tabellen und Abbildungen (Leitfaden S. 44)
% singlelinecheck=false: macht die Caption linksbündig (statt zentriert)
% labelfont auf fett: (Tabelle x.y:, Abbildung: x.y)
% font auf fett: eigentliche Bezeichnung der Abbildung oder Tabelle
% Fettschrift laut Leitfaden 2018 S. 45
\usepackage[singlelinecheck=false, labelfont=bf, font=bf]{caption}
\usepackage{caption}
\usepackage{enumitem}
\usepackage{amssymb}
\usepackage{mathptmx}
%\usepackage{minted} %Kann für schöneres Syntax Highlighting genutzt werden. ACHTUNG: Python muss installiert sein.
\usepackage[scaled=0.9]{helvet} % Behebt, zusammen mit Package courier, pixelige Überschriften. Ist, zusammen mit mathptx, dem times-Package vorzuziehen. Details: https://latex-kurs.de/fragen/schriftarten/Times_New_Roman.html
\usepackage{courier}
\usepackage{amsmath}
\usepackage[table]{xcolor}
\usepackage{marvosym}			% Verwendung von Symbolen, z.B. perfektes Eurozeichen

\renewcommand\familydefault{\sfdefault}
\usepackage{ragged2e}

% Mehrere Fussnoten nacheinander mit Komma separiert
\usepackage[hang,multiple]{footmisc}
\setlength{\footnotemargin}{2em}

% todo Aufgaben als Kommentare verfassen für verschiedene Editoren
\usepackage{todonotes}

% Verhindert, dass nur eine Zeile auf der nächsten Seite steht
\setlength{\marginparwidth}{2cm}
\usepackage[all]{nowidow}

%-----------------------------------
% Farbdefinitionen
%-----------------------------------
\definecolor{darkblack}{rgb}{0,0,0}
\definecolor{dunkelgrau}{rgb}{0.8,0.8,0.8}
\definecolor{hellgrau}{rgb}{0.0,0.7,0.99}
\definecolor{mauve}{rgb}{0.58,0,0.82}
\definecolor{dkgreen}{rgb}{0,0.6,0}

%-----------------------------------
% Pakete für Tabellen
%-----------------------------------
\usepackage{epstopdf}
\usepackage{nicefrac} % Brüche
\usepackage{multirow}
\usepackage{rotating} % vertikal schreiben
\usepackage{mdwlist}
\usepackage{tabularx}% für Breitenangabe

%-----------------------------------
% sauber formatierter Quelltext
%-----------------------------------
\usepackage{listings}
% JavaScript als Sprache definieren:
\lstdefinelanguage{JavaScript}{
	keywords={break, super, case, extends, switch, catch, finally, for, const, function, try, continue, if, typeof, debugger, var, default, in, void, delete, instanceof, while, do, new, with, else, return, yield, enum, let, await},
	keywordstyle=\color{blue}\bfseries,
	ndkeywords={class, export, boolean, throw, implements, import, this, interface, package, private, protected, public, static},
	ndkeywordstyle=\color{darkgray}\bfseries,
	identifierstyle=\color{black},
	sensitive=false,
	comment=[l]{//},
	morecomment=[s]{/*}{*/},
	commentstyle=\color{purple}\ttfamily,
	stringstyle=\color{red}\ttfamily,
	morestring=[b]',
	morestring=[b]"
}

\lstset{
	%language=JavaScript,
	numbers=left,
	numberstyle=\tiny,
	numbersep=5pt,
	breaklines=true,
	showstringspaces=false,
	frame=l ,
	xleftmargin=5pt,
	xrightmargin=5pt,
	basicstyle=\ttfamily\scriptsize,
	stepnumber=1,
	keywordstyle=\color{blue},          % keyword style
  	commentstyle=\color{dkgreen},       % comment style
  	stringstyle=\color{mauve}         % string literal style
}

%-----------------------------------
%Literaturverzeichnis Einstellungen
%-----------------------------------

% Biblatex

\usepackage{url}
\urlstyle{same}

%%%% Neuer Leitfaden IEEE
\usepackage[
backend=biber,
style=ieee,
sorting=nty, % sortierung nach name, titel , jahr 
maxcitenames=3,	% mindestens 3 Namen ausgeben bevor et. al. kommt
maxbibnames=999,
date=iso,
seconds=true, %werden nicht verwendet, so werden aber Warnungen unterdrückt.
urldate=iso,
dashed=false,
autocite=inline,
useprefix=true, % 'von' im Namen beachten (beim Anzeigen)
mincrossrefs = 1
]{biblatex}%iso dateformat für YYYY-MM-DD


%weitere Anpassungen für BibLaTex
% Opptionen für Biblatex
\ExecuteBibliographyOptions{%
  giveninits=false,
  isbn=true,
  url=true,
  doi=false,
  eprint=false,
  maxcitenames=2, % et al. ab dem 3. Autor
  mincitenames=1,
  backref=false, % Rückverweise auf Zitatseiten
  bibencoding=utf8, % wenn .bib in utf8, sonst ascii
  bibwarn=true, % Warnung bei fehlerhafter bib-Datei
}%

% et al. an Stelle von u.a.
\DefineBibliographyStrings{ngerman}{
  andothers = {{et\,al\adddot}},
}

%% et al. anstatt u. a. bei mehr als drei Autoren.
\DefineBibliographyStrings{ngerman}{
  andothers = {{et\,al\adddot}},
}
\DefineBibliographyStrings{english}{
  andothers = {{et\,al\adddot}},
}


% Klammern um das Jahr in der Fußnote
% Package biblatex Warning: Macro 'cite:labelyear+extrayear' undefined. Using \newbibmacro on input line 31.
%\renewbibmacro*{cite:labelyear+extrayear}{%
%  \iffieldundef{labelyear}
%    {}
%    {\printtext[bibhyperref]{%
%       \mkbibparens{%
%         \printfield{labelyear}%
%         \printfield{extrayear}}}}}

% Package biblatex Warning: Macro 'cite:title' undefined. Using \newbibmacro on input line 39.
%\renewbibmacro*{cite:title}{%
%  \printtext[bibhyperref]{%
%    \printfield[citetitle]{labeltitle}%
%    \setunit{\addcomma\space}%
%    \printdate}}

\DeclareNameFormat{last-first}{%
  \iffirstinits
  {\usebibmacro{name:family-given}
    {\namepartfamily}
    {\namepartgiveni}
    {\namepartprefix}
    {\namepartsuffix}
  }
  {\usebibmacro{name:family-given}
    {\namepartfamily}
    {\namepartgiven}
    {\namepartprefix}
    {\namepartsuffix}
  }%
  \usebibmacro{name:andothers}}

% Alternative Notation der Fußnoten
% Zeigt sowohl den Nachnamen als auch den Vornamen an
% Beispiel: \fullfootcite[Vgl. ][Seite 5]{Tanenbaum.2003}
\DeclareCiteCommand{\fullfootcite}[\mkbibfootnote]
{\usebibmacro{prenote}}
{\usebibmacro{citeindex}%
  \printnames[sortname][1-1]{author}%
  \addspace (\printfield{year})}
{\addsemicolon\space}
{\usebibmacro{postnote}}

%Autoren (Nachname, Vorname)
\DeclareNameAlias{default}{family-given}

%Reihenfolge von publisher, year, address verändern
% Achtung, bisher nur für den Typ @book definiert

%% Definiert @Book Eintrag
\DeclareBibliographyDriver{book}{%
  \printnames{author}%
  \newunit\addcolon\space
  \printfield{title}%
  \setunit*{,\space}%
  \printfield{edition}%
  \setunit*{\addcomma\space}%
  \printlist{publisher}%
  \newunit\newblockpunct
  \printlist{location}%
  \setunit*{\space}%
  \printfield{year}%
  \setunit*{,\space}%
  \printfield{isbn}%
  \finentry}

%% Definiert @Online Eintrag
\DeclareBibliographyDriver{online}{%
  \printnames{author}%
  \newunit\newblockpunct
  \printfield{title}%
  \setunit*{,\space}%
  %\newunit\newblock
  \printfield{url}%
  \setunit*{,\space Erscheinungsjahr:\space}%
  \printfield{year}%
  \setunit*{,\space Aufruf am:\space}%
  \printfield{note}%
  \finentry}

%% Definiert @Article Eintrag
\DeclareBibliographyDriver{article}{%
  \printnames{author}%
  \newunit\newblockpunct
  \printfield{title}%
  \setunit*{.\space In:\space}%
  %\newunit\newblock
  \usebibmacro{journal}%
  \setunit*{\space (}%
  \printfield{year}\newunit{)}%
  \finentry}

%% Definiert @InProceedings Eintrag
\DeclareBibliographyDriver{inproceedings}{%
  \printnames{author}%
  \setunit*{,\space (}%
  \printfield{year}\newunit{)}%
  \newunit\newblockpunct
  \printfield{title}%
  \setunit*{\space}%
  \usebibmacro{booktitle}%
  \setunit*{,\space}%
  \printfield{isbn}%
  \setunit*{,\space}%
  \printfield{doi}%
  \finentry}


%Aufräumen die Felder werden laut Leitfaden nicht benötigt.
\AtEveryBibitem{%
  \ifentrytype{book}{
    \clearfield{issn}%
    \clearfield{doi}%
    \clearfield{isbn}%
    \clearfield{url}
    \clearfield{eprint}
  }{}
  \ifentrytype{collection}{
    \clearfield{issn}%
    \clearfield{doi}%
    \clearfield{isbn}%
    \clearfield{url}
    \clearfield{eprint}
  }{}
  \ifentrytype{incollection}{
    \clearfield{issn}%
    \clearfield{doi}%
    \clearfield{isbn}%
    \clearfield{url}
    \clearfield{eprint}
  }{}
  \ifentrytype{article}{
    \clearfield{issn}%
    \clearfield{doi}%
    \clearfield{isbn}%
    \clearfield{url}
    \clearfield{eprint}
  }{}
  \ifentrytype{inproceedings}{
    \clearfield{issn}%
    \clearfield{doi}%
    \clearfield{isbn}%
    \clearfield{url}
    \clearfield{eprint}
  }{}
}

%Doppelpunkt nach dem letzten Autor
\renewcommand*{\labelnamepunct}{\addcolon\addspace }

%Komma an Stelle des Punktes
\renewcommand*{\newunitpunct}{\addcomma\space}

%Autoren durch Semikolon trennen
\newcommand*{\bibmultinamedelim}{\addsemicolon\space}%
\newcommand*{\bibfinalnamedelim}{\addsemicolon\space}%
\AtBeginBibliography{%
  \let\multinamedelim\bibmultinamedelim
  \let\finalnamedelim\bibfinalnamedelim
}

%Titel nicht kursiv anzeigen
\DeclareFieldFormat{title}{#1\isdot}





%%%% Ende Alter Leitfaden

%Bib-Datei einbinden
\addbibresource{literatur/literatur.bib}

% Zeilenabstand im Literaturverzeichnis ist Einzeilig
% siehe Leitfaden S. 14
\AtBeginBibliography{\singlespacing}

%-----------------------------------
% Silbentrennung
%-----------------------------------
\usepackage{hyphsubst}
\HyphSubstIfExists{ngerman-x-latest}{%
\HyphSubstLet{ngerman}{ngerman-x-latest}}{}

%-----------------------------------
% Pfad fuer Abbildungen
%-----------------------------------
\graphicspath{{./}{./abbildungen/}}

%-----------------------------------
% Weitere Ebene einfügen
%-----------------------------------
\usepackage{titletoc}

\makeatletter

% Setze die Tiefe des Inhaltsverzeichnis auf 4 Ebenen
% Damit erscheinen \paragraph-Sektionen auch im Inhaltsverzeichnis
\setcounter{secnumdepth}{4}
\setcounter{tocdepth}{4}

% Fuege Abstand nach unten wie in einer normalen \section hinzu
% Andernfalls haette \paragraph keinen Zeilenumbruch
% Der Zeilenumbruch koennte mit einer leeren \mbox{} ersetzt werden
% Jedoch klebt dann der Text relativ nah an der Ueberschrift
\renewcommand{\paragraph}{%
  \@startsection{paragraph}{4}%
  {\z@}{3.25ex \@plus 1ex \@minus .2ex}{1.5ex plus 0.2ex}%
  {\normalfont\normalsize\bfseries\sffamily}%
}

\makeatother


%-----------------------------------
% Paket für die Nutzung von Anhängen
%-----------------------------------
\usepackage{appendix}

%-----------------------------------
% Zeilenabstand 1,5-zeilig
%-----------------------------------
\usepackage{setspace}
\onehalfspacing

%-----------------------------------
% Absätze durch eine neue Zeile
%-----------------------------------
\setlength{\parindent}{0mm}
\setlength{\parskip}{0.8em plus 0.5em minus 0.3em}

\sloppy					%Abstände variieren
\pagestyle{headings}

%----------------------------------
% Präfix in das Abbildungs- und Tabellenverzeichnis aufnehmen, statt nur der Nummerierung (siehe Issue #206).
%----------------------------------
\KOMAoption{listof}{entryprefix} % Siehe KOMA-Script Doku v3.28 S.153
\BeforeStartingTOC[lof]{\renewcommand*\autodot{:}} % Für den Doppelpunkt hinter Präfix im Abbildungsverzeichnis
\BeforeStartingTOC[lot]{\renewcommand*\autodot{:}} % Für den Doppelpunkt hinter Präfix im Tabellenverzeichnis

%-----------------------------------
% Abkürzungsverzeichnis
%-----------------------------------
\usepackage[printonlyused]{acronym}

%-----------------------------------
% Symbolverzeichnis
%-----------------------------------
% Quelle: https://www.namsu.de/Extra/pakete/Listofsymbols.pdf
\usepackage[final]{listofsymbols}

%-----------------------------------
% Glossar
%-----------------------------------
\usepackage{glossaries}
\glstoctrue %Auskommentieren, damit das Glossar nicht im Inhaltsverzeichnis angezeigt wird.
\makenoidxglossaries
\newglossaryentry{glossar}{name={Glossar},description={In einem Glossar werden Fachbegriffe und Fremdwörter mit ihren Erklärungen gesammelt.}}
\newglossaryentry{glossaries}{name={Glossaries},description={Glossaries ist ein Paket was einen im Rahmen von LaTeX bei der Erstellung eines Glossar unterstützt.}}

%-----------------------------------
% PDF Meta Daten setzen
%-----------------------------------
\usepackage[hyperfootnotes=false]{hyperref} %hyperfootnotes=false deaktiviert die Verlinkung der Fußnote. Ansonsten inkompaibel zum Paket "footmisc"
% Behebt die falsche Darstellung der Lesezeichen in PDF-Dateien, welche eine Übersetzung besitzen
% siehe Issue 149
\makeatletter
\pdfstringdefDisableCommands{\let\selectlanguage\@gobble}
\makeatother

\hypersetup{
    pdfinfo={
        Title={\myTitel},
        Subject={\myStudiengang},
        Author={\myAutor},
        Build=1.1
    }
}

%-----------------------------------
% PlantUML
%-----------------------------------
%\usepackage{plantuml}

%-----------------------------------
% Umlaute in Code korrekt darstellen
% siehe auch: https://en.wikibooks.org/wiki/LaTeX/Source_Code_Listings
%-----------------------------------
\lstset{literate=
	{á}{{\'a}}1 {é}{{\'e}}1 {í}{{\'i}}1 {ó}{{\'o}}1 {ú}{{\'u}}1
	{Á}{{\'A}}1 {É}{{\'E}}1 {Í}{{\'I}}1 {Ó}{{\'O}}1 {Ú}{{\'U}}1
	{à}{{\`a}}1 {è}{{\`e}}1 {ì}{{\`i}}1 {ò}{{\`o}}1 {ù}{{\`u}}1
	{À}{{\`A}}1 {È}{{\'E}}1 {Ì}{{\`I}}1 {Ò}{{\`O}}1 {Ù}{{\`U}}1
	{ä}{{\"a}}1 {ë}{{\"e}}1 {ï}{{\"i}}1 {ö}{{\"o}}1 {ü}{{\"u}}1
	{Ä}{{\"A}}1 {Ë}{{\"E}}1 {Ï}{{\"I}}1 {Ö}{{\"O}}1 {Ü}{{\"U}}1
	{â}{{\^a}}1 {ê}{{\^e}}1 {î}{{\^i}}1 {ô}{{\^o}}1 {û}{{\^u}}1
	{Â}{{\^A}}1 {Ê}{{\^E}}1 {Î}{{\^I}}1 {Ô}{{\^O}}1 {Û}{{\^U}}1
	{œ}{{\oe}}1 {Œ}{{\OE}}1 {æ}{{\ae}}1 {Æ}{{\AE}}1 {ß}{{\ss}}1
	{ű}{{\H{u}}}1 {Ű}{{\H{U}}}1 {ő}{{\H{o}}}1 {Ő}{{\H{O}}}1
	{ç}{{\c c}}1 {Ç}{{\c C}}1 {ø}{{\o}}1 {å}{{\r a}}1 {Å}{{\r A}}1
	{€}{{\EUR}}1 {£}{{\pounds}}1 {„}{{\glqq{}}}1
}

%-----------------------------------
% Deckblatt Hintergrundbild
%-----------------------------------
\usepackage{eso-pic}
\newcommand\BackgroundPic{%
\put(0,0){%
\parbox[b][\paperheight]{\paperwidth}{%
\vfill
\centering
\includegraphics[width=\paperwidth,height=\paperheight,keepaspectratio]{abbildungen/a4_background_with_box.png}%
\vfill
}}}



%-----------------------------------
% Kopfbereich / Header definieren
%-----------------------------------
\pagestyle{fancy}
\fancyhf{}

% define font size of both header and footer 
\newcommand{\changefont}{%
    \fontsize{12}{11}\selectfont
}

% Header Seitenzahl oben, links, entsprechend Leitfaden , kein Abstand zur Linie setzen
\fancyhead[R]{\raisebox{-.5\baselineskip}[0pt][0pt]{\changefont \nouppercase\leftmark}}
\renewcommand{\headrulewidth}{1.5pt}

% Footer, set font lower 
\fancyfoot[R]{\changefont \textcolor{gray}{\myTitel} | \textcolor{teal}{\thepage}}
\renewcommand{\footrulewidth}{1.5pt}

% these rules set the header color and footer 
\renewcommand{\headrule}{\hbox to\headwidth{%
  \color{teal}\leaders\hrule height \headrulewidth\hfill}}
\renewcommand{\footrule}{\hbox to\headwidth{
    \color{teal}\leaders\hrule height \footrulewidth\hfill}}

%-----------------------------------
% Damit die hochgestellten Zahlen auch auf die Fußnote verlinkt sind (siehe Issue 169)
%-----------------------------------
\hypersetup{colorlinks=true, breaklinks=true, linkcolor=darkblack, citecolor=darkblack, menucolor=darkblack, urlcolor=darkblack, linktoc=all, bookmarksnumbered=false, pdfpagemode=UseOutlines, pdftoolbar=true}
\urlstyle{same}%gleiche Schriftart für den Link wie für den Text

%-----------------------------------
% Start the document here:
%-----------------------------------
\begin{document}

\pagenumbering{Roman}								% Seitennumerierung auf römisch umstellen
\newcolumntype{C}{>{\centering\arraybackslash}X}	% Neuer Tabellen-Spalten-Typ:
%Zentriert und umbrechbar

%-----------------------------------
% Textcommands
%-----------------------------------
\input{skripte/textcommands}

%-----------------------------------
% Deckblatt
%-----------------------------------
\begin{titlepage}
    % define background picture and title page
    \newgeometry{left=2.7cm, right=9.2cm, top=3.5cm, bottom=4cm}
    \AddToShipoutPicture*{\BackgroundPic}
    \color{teal}{
        % Information 
        \begin{huge}
            \begin{spacing}{1.125}
                \setstretch{1.0}
                \begin{flushleft}
                    \myTitel
                \end{flushleft}
            \end{spacing}
        \end{huge}
        % horizontale Linie setzen
        \par\noindent\rule{\textwidth}{2pt}

        % datum und name getrennt durch spacing mit Authorenname
        \vspace{40em}

        \par\noindent\rule{\textwidth}{2pt}
        % email adresse 
        \myVorname.\myNachname@fom-net.de \break
        \today

    }
\end{titlepage}

%-----------------------------------
% Titlepage
%-----------------------------------
\begin{titlepage}
	\newgeometry{left=2cm, right=2cm, top=2cm, bottom=2cm}
	\begin{center}
		\includegraphics[width=10cm]{abbildungen/FOM_Logo.png} \\ % FOM Logo und text 
		\vspace{.5cm}
		\begin{Large}\langde{Hochschulzentrum}\langen{university location} \myHochschulStandort\end{Large}\\
		\vspace{2cm}
		\begin{Large}\textbf{\myThesisArt}\end{Large}\\
		\vspace{.5cm}
		% \langde{Berufsbegleitender Studiengang}
		% \langen{part-time degree program}\\
		% \mySemesterZahl. Semester\\
		\langde{im Studiengang}\langen{in the study course} \myStudiengang
		\vspace{1.7cm}

		\langde{zur Erlangung des Grades eines}\langen{to obtain the degree of}\\
		\vspace{0.5cm}
		\begin{Large}{\myAkademischerGrad}\end{Large}\\
		% Oder für Hausarbeiten:
		%\textbf{im Rahmen der Lehrveranstaltung}\\
		%\textbf{\myLehrveranstaltung}\\
		\vspace{1.8cm}
		\langde{über das Thema}
		\langen{on the subject}\\
		\vspace{0.5cm}
		\large{\textbf{\myTitel}}\\
		\vspace{2cm}
		\langde{von}\langen{by}\\
		\vspace{0.5cm}
		\begin{Large}{\myAutor}\end{Large}\\
	\end{center}
	\normalsize
	\vfill
	\begin{tabular}{ l l }
		\langde{Betreuer} % für Hausarbeiten
		%\langde{Erstgutachter} % für Bachelor- / Master-Thesis
		\langen{Advisor}:              & \myBetreuer    \\
		\langde{Matrikelnummer}
		\langen{Matriculation Number}: & \myMatrikelNr  \\
		\langde{Abgabedatum}
		\langen{Submission}:           & \myAbgabeDatum
		\\
	\end{tabular}
\end{titlepage}


%-----------------------------------
% Inhaltsverzeichnis
%-----------------------------------
% Um das Tabellen- und Abbbildungsverzeichnis zu de/aktivieren ganz oben in Documentclass schauen
\setcounter{page}{2}
\addtocontents{toc}{\protect\enlargethispage{-20mm}}% Die Zeile sorgt dafür, dass das Inhaltsverzeichnisseite auf die zweite Seite gestreckt wird und somit schick aussieht. Das sollte eigentlich automatisch funktionieren. Wer rausfindet wie, kann das gern ändern.
\setcounter{tocdepth}{4}
\tableofcontents
\newpage

%-----------------------------------
% Abbildungsverzeichnis
%-----------------------------------
\listoffigures
\newpage
%-----------------------------------
% Tabellenverzeichnis
%-----------------------------------
\listoftables
\newpage
%-----------------------------------
% Abkürzungsverzeichnis
%-----------------------------------
% Falls das Abkürzungsverzeichnis nicht im Inhaltsverzeichnis angezeigt werden soll
% dann folgende Zeile auskommentieren.
\addcontentsline{toc}{section}{\abbreHeadingName}
\section*{\langde{Abkürzungsverzeichnis}\langen{List of Abbreviations}}

\begin{acronym}[WYSIWYG]\itemsep0pt %der Parameter in Klammern sollte die längste Abkürzung sein. Damit wird der Abstand zwischen Abkürzung und Übersetzung festgelegt
  \acro{OC}{FOM Online Campus}
  \acro{WYSIWYG}{What you see is what you get}
  \acro{Beispiel}{Nicht verwendet, taucht nicht im Abkürzungsverzeichnis auf}
\end{acronym}
\newpage

%-----------------------------------
% Symbolverzeichnis
%-----------------------------------
% In Overleaf führt der Einsatz des Symbolverzeichnisses zu einem Fehler, der aber ignoriert werdne kann
% Falls das Symbolverzeichnis nicht im Inhaltsverzeichnis angezeigt werden soll
% dann folgende Zeile auskommentieren.
%\addcontentsline{toc}{section}{\symheadingname}
%\input{skripte/symbolDef}
%\listofsymbols
%\newpage

%-----------------------------------
% Glossar
%-----------------------------------
\printnoidxglossaries
\newpage

%-----------------------------------
% Sperrvermerk
%-----------------------------------
\input{kapitel/anhang/sperrvermerk}

%-----------------------------------
% Seitennummerierung auf arabisch und ab 1 beginnend umstellen
%-----------------------------------
\pagenumbering{arabic}
\setcounter{page}{1}

%-----------------------------------
% Kapitel / Inhalte
%-----------------------------------
% Hinzugefügt aufgrund von Issue 167
%-----------------------------------
% Kapitel / Inhalte
%-----------------------------------
\input{kapitel/einleitung/einleitung}
\input{kapitel/kapitel_1/kapitel_1}
\newpage
\section{Latex-Details} \label{latexDetails}

\subsection{Verwendete Software, Editor und Zusatzpakete}
\subsubsection{Windows 8+}
\begin{itemize}
	\item MikTex: 2.9, 32-bit
	\item Biblatex: 3.5, Zusatz: Biber.exe
	\item Editor: TexStudio (kann ich empfehlen), Notepad++
\end{itemize}

\subsubsection{Mac OSX und iOS}
\begin{itemize}
	\item MacTeX: \url{https://tug.org/mactex}
	\item Editor: TexPad \url{https://www.texpadapp.com}
\end{itemize}

\subsubsection{Online}
Overleaf ist eine Online-Anwendung mit der Ihr direkt im Browser an eurer Thesis schreiben könnt. Bis 1GB Größe und maximal 60 Einzeldateien könnt ihr Overleaf kostenlos nutzen: \url{https://www.overleaf.com/}


\subsection{Dokumentenklasse}
Eigentlich hatte Prof. Finke empfohlen die Dokumentklassen \enquote{Book} oder \enquote{Report} für die Erstellung der Bachelor-Thesis zu verwenden, da diese über weitere Gliederungsebenen verfügen. Ich verwende dennoch eine leicht modifizierte Komaskript-Klasse \enquote{scrartcl}, mit der Erweiterung um eine Ebene. Siehe (skripte/weitereEbene.tex). Das Skript stammt irgendwo aus den Netz und übersteigt meine \LaTeX{}-Fähigkeiten. Dadurch kann ich über eine weitere Ebene in der Arbeit verfügen, ohne mich mit der Modifikation von Kapitel-Seiten rumschlagen~\footcite[Vgl. ][S. 5]{Tanenbaum.2003} zu müssen. Diese Quelle ist nur zur Demonstration und hat keinen inhaltlichen Bezug hierzu. Es werden übrigens nur die Quellen im Literaturverzeichnis angezeigt, die auch referenziert sind.


\subsection{Grafiken}
Das Paket \textbackslash usepackage\{float\} ermöglicht es die Grafiken und Tabellen an der Stelle im Text zu positionieren, wo diese im Quelltext stehen (Option H). Ansonsten würde \LaTeX{} diese dort unterbringen, wo es typographisch sinnvoll wäre - das wollen wir ja nicht ;-).

Die Breite der Grafiken am Besten relativ zum Text angeben.

\subsection{Quellcode}
Quellcode kann auf unterschiedliche Arten eingebaut werden.
Zum einen kann es hier durch direktives Einbinden in der Kapitel-Datei geschehen.
\begin{lstlisting}
% Hier wird aufgezeigt, wie man eine Grafik einbindet, es wird also in der PDF angezeigt,
% da es in einem Quellcode-Listing steht.
% Auch wenn es hier faelschlicherweise als LaTeX-Befehl angezeigt wird.
\includegraphics[width=0.9\textwidth]{sup}
\end{lstlisting}

Bei längeren Quellcode-Listings empfiehlt es sich jedoch auf eine externe Datei im Ordner Quellcode zu verlinken und diese einzubauen:
\lstinputlisting[language=HTML]{./Quellcode/Beispiel.html}

Statt dem Package lstlisting, welches direkt auf Tex basiert, kann auch das Package minted verwendet werden.
Dieses Package basiert auf python-pygments und unterstützt weit mehr Sprachkonstrukte als lstlisting.
Um das Paket zu verwenden muss es eingebunden werden und zusätzlich python-pygments installiert sein.
(Dies ist mit im Dockerfile vorhanden. Für die anderen Compile-Methoden, wie das native verwenden von Tex Live findet sich hier die Installationsanleitung für das minted Paket: https://ctan.org/pkg/minted?lang=de)

Damit das kompilieren ohne Python trotzdem möglich ist, ist die Funktion standardmäßig ausgebaut. Deshalb muss zusätzlich in der Datei \begin{verbatim}thesis_main.tex \usepackage{minted} \end{verbatim} wieder einkommentiert werden.

Minted lässt sich dann ganz ähnlich zu lstlisting verwenden:
\begin{lstlisting}
	\begin{minted}{c}
		int main() {
			printf("hello, world");
			return 0;
		}
	\end{minted}
\end{lstlisting}

Da der Pfad zu den Abbildungen im Hauptdokument definiert wurde, muss hier nur noch der Name des Bildes ohne Dateiendung stehen (sup).

\begin{figure}[H]
	\caption{Titel der Abbildung hier}
	\includegraphics[width=0.9\textwidth]{sup}
	\\
	Quelle: Eigene Darstellung
\end{figure}

\subsection{Tabellen}
\begin{table}[H]
	\caption{Beispieltabelle 1}
	\label{tbl:beispieltabelle2}
	\begin{tabularx}{\textwidth}[ht]{|l|X|l|}
		\hline
		\textbf{Abkürzung} & \textbf{Beschreibung}             & \textbf{Berechnung} \\
		\hline\hline
		MEK                & Materialeinzelkosten              &                     \\
		MGK                & Materialgemeinkosten              & $+ \uparrow$~*      \\
		FEK                & Fertigungseinzelkosten            &                     \\
		FGK                & Fertigungsgemeinkosten            & $+ \uparrow$~*      \\
		SEKF               & Sondereinzelkosten der Fertigung  &                     \\
		\hline\hline
		\multicolumn{3}{|l|}{\textbf{= Herstellungskosten}}                          \\
		\hline\hline
		VwGK               & Verwaltungsgemeinkosten           & $+ \uparrow$~*      \\
		VtGK               & Vertriebsgemeinkosten             & $+ \uparrow$~*      \\
		SEKVt              & Sondereinzelkosten des Vertriebes &                     \\
		\hline\hline
		\multicolumn{3}{|l|}{\textbf{= Selbstkosten}}                                \\
		\hline\hline
		\multicolumn{3}{|l|}{+ Gewinnaufschlag}                                      \\
		\multicolumn{3}{|l|}{+ Rabatte}                                              \\
		\hline\hline
		\multicolumn{3}{|l|}{\textbf{= Nettoverkaufspreis (NVP)}}                    \\
		\hline
		\multicolumn{3}{|l|}{+ Umsatzsteuer}                                         \\
		\hline\hline
		\multicolumn{3}{|l|}{\textbf{= Bruttoverkaufspreis (BVP)}}                   \\
		\hline
	\end{tabularx} \\
	\cite[Quelle: In Anlehnung an][S. 4]{Beckert.2012}
\end{table}

%\clearpage % hiermit werden alle Bilder Tabellen ausgeworfen

\subsection{Biblatex}
\subsubsection{Erklärung}
Von den vielen verfügbaren Literatur-Paketen habe ich mich für Biblatex entschieden. Die Anforderungen der FOM sollten hiermit erfüllt sein. Ich habe bisher nur Einträge \enquote{@book} getestet. Wie immer steckt der Teufel hier im Detail und es wird sich später herausstellen, ob Biblatex eine gute Wahl war. Die Anpassungen hierfür liegen unter skripte/modsBiblatex. Ich verwende das Backend Biber, welches bib-Dateien in UTF-8 verarbeiten kann.

In der für den Leitfaden 2018 aktualisierten Version sind außerdem Beispiele für \enquote{online},\footcite[Vgl.][]{website:angular:aboutAngular} also Webseiten, und \enquote{article},\footcite[Vgl.][S. 140]{Decker2009} also wissenschaftliche Artikel, enthalten.

Laut Leitfaden sollen maximal 3 Autoren genannt werden und danach mit
\enquote{et. al.} bzw. \enquote{u.a.} ergänzt werden. Damit im Literaturverzeichnis auch nur max.
3 Autoren stehen, muss man beim Füllen der literatur.bib-Datei darauf achten auch nur 3
einzutragen. Weitere Autoren kann man einfach mit \enquote{and others} ergänzen.
Siehe Eintrag für \enquote{Balzert.2008}. Zitiert man dann diese Werk, werden auch in
der Fussnote alle Autoren korrekt genannt wie in dieser
Fußnote\footcite[Vgl.][S. 1]{Balzert.2008} zu sehen ist.

Hat man dagegen mehr als 3 Autoren in der bib-Datei hinterlegt, stehen im
Literaturverzeichnis alle drin. In der Fussnote dagegen, steht nur
einer\footcite[Vgl.][S. 1]{Balzert2.2008}, was dem Leitfaden widerspricht.

Die Anzahl von 3 wird übrigens über die Option \enquote{maxcitenames=3} des
biblatex-Packages gesetzt. Man muss selbst schauen, dass die Anzahl der Autoren
in den Bib-Dateien mit der Optionseinstellung übereinstimmt.

\subsubsection{Beispielfußnoten}
Diese Fussnote soll zeigen, wie mit einem \enquote{von} vor dem Namen des Autors
umgegangen wird\footcite[Vgl.][S. 1]{Lucke2018}. Man muss für die korrekte
Sortierung eines solchens Namens im Literaturverzeichnis einen \enquote{sortkey}
setzen.

Diese Fussnote soll zeigen, wie mit einer Online-Quelle ohne Jahresangabe
umgegangen wird\footcite[Vgl.][]{Belastingdienst}.

Diese Fußnote\footcite[Vgl.][S.1]{Beckert.2012} ist nur dazu da zu zeigen, wie mit mehreren Quellen des selben Autors aus dem selben Jahr umgegangen wird, wenn das Stichwort gleich bleibt \footcite[Vgl.][S.2]{Beckert.2012.1} oder sich ändert\footcite[Vgl.][S.3]{Beckert.2012.2}. Laut Leitfaden sollte bei gleichem Autor, Jahr und Stichwort ein Buchstabe an die Jahreszahl gehangen werden. Zum Beispiel 2012a.

Die folgenden Fußnoten dienen dazu zu zeigen, dass die Nummern von zwei direkt aufeinanderfolgende Fußnoten mit Komma getrennt werden.\footcite[Vgl.][S.2]{Beckert.2012.1}\footcite[Vgl.][S. 1]{Lucke2018}
\subsection{Abkürzungen}
Abkürzungen werden mithilfe des Pakets Acronym eingebunden. Alle Abkürzungen sollten in der Datei acronyms.tex mithilfe des \begin{verbatim}
	\acro
\end{verbatim} Befehls festgelegt werden. Im Text werden diese dann mit \begin{verbatim}
	\ac{Abkürzung}
\end{verbatim} benutzt. Bei der ersten Verwendung einer Abkürzung wird der Begriff in beiden Formen dargestellt. So wie hier: \ac{WYSIWYG}. Nur wenn eine Abkürzung tatsächlich verwendet wird erscheint sie auch im Abkürzungsverzeichnis.

Sollte es im Abkürzungsverzeichnis zu Anzeigefehlern kommen kann dies daher rühren, dass eine Abkürzung verwendet wird, die länger ist als \ac{WYSIWYG}. In diesem Fall müsst ihr in der Datei acronyms.tex den Parameter [WYSIWYG] durch eure längere Abkürzung ersetzen.

\subsection{Formeln}
Um eine Formel nach links aus zurichten muss sie zwischen \& und \& eingesetzt werden:

\textbf{Formel 1: Erste Formel}
\begin{flalign}
	& L_P{=} 10lg \cdot \frac{P}{1 mW} &
\end{flalign}
\cite[Quelle: In Anlehnung an][S. 4]{Beckert.2012}


Etwas mehr Text.

Ansonsten wird sie mittig ausgerichtet test.
% Mehr infos: http://www.ctex.org/documents/packages/math/amsldoc.pdf

\textbf{Formel 2: Zweite Formel}
\begin{flalign}
	L_P{=} 10lg \cdot \frac{P}{1 mW}
\end{flalign}
\cite[Quelle: In Anlehnung an][S. 4]{Beckert.2012}

\subsection{Symbole}
% die folgenden Symbole haben nicht mit der Formel oben drüber zu tun
% Das hier ist ein definiertes Symbol: \symnz und das hier auch \AB . Symbole werden in der Datei Skripte symboldef.tex zentral definiert.

\subsection{Glossar}
Begriffserklärungen bzw. das \gls{glossar} wird mithilfe des Pakets \gls{glossaries} eingebunden. Alle Begriffe die erklärt werden sollen, sollten in der Datei glossar.tex mithilfe des \begin{verbatim}
	\newglossaryentry
\end{verbatim} Befehls festgelegt werden. Im Text werden diese dann mit \begin{verbatim}
	\gls{Begriff}
\end{verbatim} benutzt.


\subsection{Listen und Aufzählungen}
\subsubsection{Listen}
\begin{itemize}
	\item ein wichtiger Punkt
	\item noch ein wichtiger Punkt
	\item und so weiter
\end{itemize}
\subsubsection{Aufzählungen}
\begin{enumerate}
	\item Reihenfolge ist hier wichtig
	\item Dieser Punkt kommt nach dem ersten
	\item Da sollte jetzt eine 3 vorne stehen
\end{enumerate}

\paragraph{Tiefste Ebene 1}
Dies ist die tiefste Gliederungsebene. Sollten doch mehr Ebenen benötigt werden, muss eine andere Dokumentenklasse verwendet werden.

\paragraph{Tiefste Ebene 2}
Der zweite Punkt in dieser Ebene ist zur Erinnerung daran, dass es nie nie niemals nur einen Unterpunkt geben darf.

\subsection{Skript zum Kompilieren}
Latex will ja bekanntlich in einer bestimmten Reihenfolge aufgerufen werden:
\begin{lstlisting}
lualatex thesis_main.tex
biber thesis_main
lualatex thesis_main.tex
lualatex thesis_main.tex
thesis_main.pdf
\end{lstlisting}

Dies ist der Inhalt der Batchdatei \enquote{compile.bat}.

\subsection{PlantUML}

\begin{lstlisting}
\begin{plantuml}
@startuml
Class01 <|-- Class02
Class03 *-- Class04
Class05 o-- Class06
Class07 .. Class08
Class09 -- Class10
@enduml
\end{plantuml}
\end{lstlisting}

\input{kapitel/fazit/fazit}

%-----------------------------------
% Apendix / Anhang
%-----------------------------------
\newpage
\section*{\AppendixName} %Überschrift "Anhang", ohne Nummerierung
\addcontentsline{toc}{section}{\AppendixName} %Den Anhang ohne Nummer zum Inhaltsverzeichnis hinzufügen

\begin{appendices}
	% Nachfolgende Änderungen erfolgten aufgrund von Issue 163
	\makeatletter
	\renewcommand\@seccntformat[1]{\csname the#1\endcsname:\quad}
	\makeatother
	\addtocontents{toc}{\protect\setcounter{tocdepth}{0}} %
	\renewcommand{\thesection}{\AppendixName\ \arabic{section}}
	\renewcommand\thesubsection{\AppendixName\ \arabic{section}.\arabic{subsection}}
	\section{Beispielanhang}\label{Beispielanhang}
Dieser Abschnitt dient nur dazu zu demonstrieren, wie ein Anhang aufgebaut seien kann.
\subsection{Weitere Gliederungsebene}
Auch eine zweite Gliederungsebene ist möglich.
\section{Bilder}
Auch mit Bildern.
Diese tauchen nicht im Abbildungsverzeichnis auf.
\begin{figure}[H]
	\centering
	\caption[]{Beispielbild}
	\label{fig:Beispielbild}
	\includegraphics[width=1\textwidth]{verzeichnisStruktur}
\end{figure}
\end{appendices}
\addtocontents{toc}{\protect\setcounter{tocdepth}{2}}

%-----------------------------------
% Literaturverzeichnis
%-----------------------------------
\newpage

% Die folgende Zeile trägt ALLE Werke aus literatur.bib in das
% Literaturverzeichnis ein, egal ob sie zietiert wurden oder nicht.
% Der Befehl ist also nur zum Test der Skripte sinnvoll und muss bei echten
% Arbeiten entfernt werden.
%\nocite{*}

%\addcontentsline{toc}{section}{Literatur}

% Die folgenden beiden Befehle würden ab dem Literaturverzeichnis wieder eine
% römische Seitennummerierung nutzen.
% Das ist nach dem Leitfaden nicht zu tun. Dort steht nur dass 'sämtliche
% Verzeichnisse VOR dem Textteil' römisch zu nummerieren sind. (vgl. S. 3)
%\pagenumbering{Roman} %Zähler wieder römisch ausgeben
%\setcounter{page}{4}  %Zähler manuell hochsetzen

% Ausgabe des Literaturverzeichnisses

% Keine Trennung der Werke im Literaturverzeichnis nach ihrer Art
% (Online/nicht-Online)
%\begin{RaggedRight}
%\printbibliography
%\end{RaggedRight}

% Alternative Darstellung, die laut Leitfaden genutzt werden sollte.
% Dazu die Zeilen auskommentieren und folgenden code verwenden:

% Literaturverzeichnis getrennt nach Nicht-Online-Werken und Online-Werken
% (Internetquellen).
% Die Option nottype=online nimmt alles, was kein Online-Werk ist.
% Die Option heading=bibintoc sorgt dafür, dass das Literaturverzeichnis im
% Inhaltsverzeichnis steht.
% Es ist übrigens auch möglich mehrere type- bzw. nottype-Optionen anzugeben, um
% noch weitere Arten von Zusammenfassungen eines Literaturverzeichnisse zu
% erzeugen.
% Beispiel: [type=book,type=article]
\printbibliography[nottype=online,heading=bibintoc,title={\langde{Literaturverzeichnis}\langen{Bibliography}}]

% neue Seite für Internetquellen-Verzeichnis
\newpage

% Laut Leitfaden 2018, S. 14, Fussnote 44 stehen die Internetquellen NICHT im
% Inhaltsverzeichnis, sondern gehören zum Literaturverzeichnis.
% Die Option heading=bibintoc würde die Internetquelle als eigenen Eintrag im
% Inhaltsverzeicnis anzeigen.
%\printbibliography[type=online,heading=bibintoc,title={\headingNameInternetSources}]
\printbibliography[type=online,heading=subbibliography,title={\headingNameInternetSources}]

\newpage
\thispagestyle{empty}

%-----------------------------------
% Ehrenwörtliche Erklärung
%-----------------------------------
\section*{%
  \langde{Ehrenwörtliche Erklärung}
  \langen{Declaration in lieu of oath}}
\langde{Hiermit versichere ich, dass die vorliegende Arbeit von mir selbstständig und ohne unerlaubte Hilfe angefertigt worden ist, insbesondere dass ich alle Stellen, die wörtlich oder annähernd wörtlich aus Veröffentlichungen entnommen sind, durch Zitate als solche gekennzeichnet habe. Ich versichere auch, dass die von mir eingereichte schriftliche Version mit der digitalen Version übereinstimmt. Weiterhin erkläre ich, dass die Arbeit in gleicher oder ähnlicher Form noch keiner Prüfungsbehörde/Prüfungsstelle vorgelegen hat. Ich erkläre mich damit \textcolor{red}{einverstanden/nicht einverstanden}, dass die Arbeit der Öffentlichkeit zugänglich gemacht wird. Ich erkläre mich damit einverstanden, dass die Digitalversion dieser Arbeit zwecks Plagiatsprüfung auf die Server externer Anbieter hochgeladen werden darf. Die Plagiatsprüfung stellt keine Zurverfügungstellung für die Öffentlichkeit dar.}
\langen{I hereby declare that I produced the submitted paper with no assistance from any other party and without the use of any unauthorized aids and, in particular, that I have marked as quotations all passages which are reproduced verbatim or near-verbatim from publications. Also, I declare that the submitted print version of this thesis is identical with its digital version. Further, I declare that this thesis has never been submitted before to any examination board in either its present form or in any other similar version. I herewith \textcolor{red}{agree/disagree} that this thesis may be published. I herewith consent that this thesis may be uploaded to the server of external contractors for the purpose of submitting it to the contractors’ plagiarism detection systems. Uploading this thesis for the purpose of submitting it to plagiarism detection systems is not a form of publication.}


\par\medskip
\par\medskip

\vspace{5cm}

\begin{table}[H]
	\centering
	\begin{tabular*}{\textwidth}{c @{\extracolsep{\fill}} ccccc}
		\myOrt, \the\day.\the\month.\the\year
		&
		% Hinterlege deine eingescannte Unterschrift im Verzeichnis /abbildungen und nenne sie unterschrift.png
		% Bilder mit transparentem Hintergrund können teils zu Problemen führen
		\includegraphics[width=0.35\textwidth]{unterschrift}\vspace*{-0.35cm}
		\\
		\rule[0.5ex]{12em}{0.55pt} & \rule[0.5ex]{12em}{0.55pt} \\
		\langde{(Ort, Datum)}\langen{(Location, Date)} & \langde{(Eigenhändige Unterschrift)}\langen{(handwritten signature)}
		\\
	\end{tabular*} \\
\end{table}

\end{document}
